% !TeX program = xelatex
% Standalone revision of main.tex; data cutoff remains 2026-10-17.
% Evidence reference: myx / 5d349e82. Revision checked at 5db371f7.
% Original main.tex is preserved. No new experiment or VM change.
\documentclass[UTF8,aspectratio=169,10pt,fontset=fandol]{ctexbeamer}
\usepackage{booktabs}
\usepackage{listings}
\setmonofont{Latin Modern Mono}

\usetheme{default}
\usefonttheme{professionalfonts}
\setbeamertemplate{navigation symbols}{}
\setbeamertemplate{headline}{}
\setbeamertemplate{itemize item}{\raisebox{0.15ex}{\scriptsize$\bullet$}}
\setbeamercolor{normal text}{fg=black,bg=white}
\setbeamercolor{structure}{fg=black}
\setbeamercolor{frametitle}{fg=black,bg=white}
\setbeamercolor{footline}{fg=black,bg=white}
\setbeamerfont{frametitle}{size=\large,series=\bfseries}
\setbeamersize{text margin left=8mm,text margin right=8mm}
\setlength{\parskip}{1pt}

% detokenize deliberately accepts literal underscores in this macro's argument.
\newcommand{\code}[1]{\texttt{\detokenize{#1}}}
\newcommand{\PerfNote}{内部单调计时包含首次调用；无计时前 workload 预跑。}
\newcommand{\DiagNote}{诊断数据，不作正式性能结果。}
\newcommand{\PageSource}{}
\newcommand{\PageMode}{}
\newcommand{\PageMeta}[2]{\gdef\PageSource{#1}\gdef\PageMode{#2}}
% No footer: environment, flags and evidence references stay in notes.md.
\setbeamertemplate{footline}{}

\lstdefinestyle{codeblock}{
  basicstyle=\ttfamily\fontsize{9.4}{10.6}\selectfont,
  keywordstyle=\color{black},commentstyle=\color{black},stringstyle=\color{black},
  columns=fullflexible,keepspaces=true,showstringspaces=false,
  breaklines=true,breakatwhitespace=true,
  postbreak=\mbox{\scriptsize$\hookrightarrow$\space},
  aboveskip=2pt,belowskip=2pt,frame=none,tabsize=2}
\lstset{style=codeblock}
\title{QuickJS与V8：字节码和属性访问分析}
\author{马云霄}
\date{}

\begin{document}

% -------------------- Main 1/6 --------------------
\PageMeta{S1：numeric\_loop 源码、真实 dump、summary.csv。}{\PerfNote}
\begin{frame}[fragile,t]{字节码和属性访问分析}
\smallskip
\textbf{数值求和：相同 JavaScript}
\begin{columns}[T,onlytextwidth]
  \begin{column}{0.42\textwidth}
\begin{lstlisting}
function wbNumericLoop(n) {
    var sum = 0;
    for (var i = 0; i < n; i++) {
        sum += i;
    }
    return sum;
}
\end{lstlisting}
    输入 \code{n=4096}；\par
    返回8386560，未溢出int32。
  \end{column}
  \begin{column}{0.18\textwidth}
    \textbf{QuickJS 摘录}
\begin{lstlisting}
get_loc1 1: i
add_loc 0: sum
inc_loc 1: i
goto8 4
\end{lstlisting}
    \code{add_loc} 更新sum\par
    并弹栈；\par
    \code{inc_loc} 更新i。
  \end{column}
  \begin{column}{0.37\textwidth}
    \textbf{Ignition 摘录}
\begin{lstlisting}
Ldar r1
Add r0, EmbeddedFeedback[None]
Mov r0, r2
Star0
Ldar r1
Inc EmbeddedFeedback[None]
Star1
JumpLoop [21], [0], FBV[0]
\end{lstlisting}
    \code{r0=sum}，\code{r1=i}；\par
    \code{Mov} 复制旧sum，\code{Star0}写新sum。
  \end{column}
\end{columns}
\smallskip
\textbf{Median（$\mu$s/完整调用）：}
QJS 43.782；V8 54.839；V8/QJS 1.253。
\par 每完整调用含4096轮；每引擎30个新进程，点估计。
\par \textbf{本例QJS更快}；虚拟槽不是CPU寄存器，指令条数/字节数不等于成本。
\end{frame}

% -------------------- Main 2/6 --------------------
\PageMeta{S2：cases/summary.csv；本页inline含obj.bias，与第1页求和不同。}{\PerfNote}
\begin{frame}[fragile,t]{四组调用方式}
\small 共同计算：\code{i=0..4095}、\code{bias=3}，累加\code{i+bias}，结果8398848。
\smallskip
\centering
\setlength{\tabcolsep}{4pt}
\renewcommand{\arraystretch}{1.17}
\begin{tabular}{@{}llrrr@{}}
  \toprule
  组别 & 循环中的主要表达式 & QJS / $\mu$s & V8 / $\mu$s & V8/QJS \\
  \midrule
  直接计算 & \code{sum += i + obj.bias} & 67.132 & 60.388 & 0.900 \\
  普通调用 & \code{sum += stepPlain(obj,i)} & 159.969 & 116.395 & 0.728 \\
  own 方法 & \code{sum += obj.step(i)} & 217.156 & 120.905 & \textbf{0.557} \\
  prototype 方法 & \code{sum += obj.step(i)} & 225.086 & 125.348 & \textbf{0.557} \\
  \bottomrule
\end{tabular}
\par\smallskip
\raggedright
单位：$\mu$s/完整调用（4096轮，含摊薄driver）；每组30进程/引擎，共240样本。
\smallskip
\begin{columns}[T,onlytextwidth]
  \begin{column}{0.44\textwidth}
    \textbf{普通函数 / 方法}
\begin{lstlisting}
function stepPlain(obj, x) {
    return x + obj.bias;
}
function stepMethod(x) {
    return x + this.bias;
}
\end{lstlisting}
  \end{column}
  \begin{column}{0.53\textwidth}
    \textbf{own与prototype：相同的循环、方法、}\par
    \textbf{参数传递与this绑定形式。}\par
    主要改变计时外的对象建立方式、\par
    属性所在位置及相关布局/反馈。\par\smallskip
    \textbf{普通函数与方法}的对比，才另外涉及\par
    显式obj参数与receiver传递差异。
  \end{column}
\end{columns}
\smallskip
一层原型不足以解释差距；\textbf{组间时间不能直接相减得到单独操作成本。}
\end{frame}

% -------------------- Main 3/6 --------------------
\PageMeta{S3：原 richards.js；site\_samples.csv、sampling\_summary.json。}{\DiagNote}
\begin{frame}[fragile,t]{原始 Richards}
\small Richards是模拟任务调度的JavaScript基准；原算法、方法读取和断言均未改。
\medskip
\begin{columns}[T,onlytextwidth]
  \begin{column}{0.54\textwidth}
    \textbf{Scheduler.schedule 摘录}\par
    此处\code{this}是Scheduler。
\begin{lstlisting}
this.currentTcb = this.currentTcb.run();
\end{lstlisting}
    \smallskip
    \textbf{TaskControlBlock.run 摘录}\par
    进入后，\code{this}是当前TCB。
\begin{lstlisting}
return this.task.run(packet);
\end{lstlisting}
    \smallskip
    \code{currentTcb}指向TCB；\par
    TCB的\code{task}指向具体任务。\par
    取得对象、读方法、调用函数分开看。
  \end{column}
  \begin{column}{0.43\textwidth}
    \textbf{按属性聚合的查找区域样本}
    \smallskip
    \centering
    \begin{tabular}{@{}lrr@{}}
      \toprule
      属性 & 样本 & 占1229 \\
      \midrule
      \code{currentTcb} & 280 & 22.8\% \\
      \code{state} & 204 & 16.6\% \\
      两处\code{run} & 79 & 6.4\% \\
      \bottomrule
    \end{tabular}
    \par\smallskip
    \raggedright
    3个诊断进程，共5393个内部样本；\par
    1229个查找区域样本已映射。\par\smallskip
    同名属性包含多个读/写位置；\par
    \textbf{分母不是全程序耗时。}
  \end{column}
\end{columns}
\medskip
\textbf{普通状态/调度字段值得继续看；目前选出的是候选位置。}
\par\smallskip
189个活跃opcode位置；N1/N2计数倍增核对，计数与采样分开。
\end{frame}

% -------------------- Main 4/6 --------------------
\PageMeta{S3：原函数最终dump；栈语义见 quickjs.c:18023--18040、18881--18901。}{\DiagNote}
\begin{frame}[fragile,t]{选定 this.state：零比较与短路怎样执行}
\begin{lstlisting}
return (this.state & STATE_HELD) != 0 || (this.state == STATE_SUSPENDED);
\end{lstlisting}
\small 调用者\code{schedule}的\code{this}是Scheduler；进入后是当前TCB；\code{r0}为虚拟槽。
\begin{columns}[T,onlytextwidth]
  \begin{column}{0.38\textwidth}
    \textbf{QuickJS：offset 2--19}
\begin{lstlisting}
 2 get_loc0 0: this
 3 get_field state
 8 get_var 0: STATE_HELD
11 and
12 push_0 0
13 neq
14 dup
15 if_true8 28
17 drop
18 get_loc0 0: this
19 get_field state
\end{lstlisting}
  \end{column}
  \begin{column}{0.59\textwidth}
    \textbf{Ignition：offset 0--19}
\begin{lstlisting}
 0 GetNamedProperty <this>, [0:"state"], FBV[0]
 4 Star0
 5 LdaGlobal [1:"STATE_HELD"], FBV[2]
 8 BitwiseAnd r0, EmbeddedFeedback[None]
11 Star0
12 LdaZero
13 TestEqual r0, EmbeddedFeedback[None]
16 LogicalNot
17 JumpIfTrue [13] (0x1eef0100166e @ 30)
19 GetNamedProperty <this>, [0:"state"], FBV[4]
\end{lstlisting}
  \end{column}
\end{columns}
\small QJS栈底$\to$栈顶：\code{[t,0]}经\code{neq}得\code{[b]}，\code{dup}得\code{[b,b]}；t为按位与结果。\par
b为左侧布尔值。分支弹一份b：真跳28返回；假经\code{drop}丢弃左值，再求右侧。\par
\textbf{片段省略处：}QJS 0--1/24--28，V8 23--30；完整序列见附录。
\end{frame}

% -------------------- Main 5/6 --------------------
\PageMeta{QJS 04be246，quickjs.c:19126--19133；S3计数与同进程IC日志。}{\DiagNote}
\begin{frame}[fragile,t]{这一处获取信息}
\small
\begin{columns}[T,onlytextwidth]
  \begin{column}{0.61\textwidth}
    \textbf{QJS：首次state读取，offset 3}\par
    每Run\textbf{10671次}；均own数据属性、整数结果；\par
    每次一次自身探查，原型命中/fallback均为0。
\begin{lstlisting}
prs = find_own_property(&pr, p, atom); \
if (prs) { \
    /* found */ \
    if (unlikely(prs->flags & JS_PROP_TMASK)) \
            goto name ## _slow_path; \
    val = JS_DupValue(ctx, pr->u.value); \
    break; \
} \
\end{lstlisting}
    \code{GET_FIELD_INLINE}内部片段，非完整快路径。\par
    省略对象检查、未命中/原型/慢路径和最终写回。\par
    receiver后续引用处理（19157--19158）见讲稿。
  \end{column}
  \begin{column}{0.36\textwidth}
    \textbf{V8：同进程IC日志}
    \smallskip
    \centering
    \setlength{\tabcolsep}{3pt}
    \begin{tabular}{@{}ccc@{}}
      \toprule
      offset & 反馈槽 & 转换码 \\
      \midrule
      0 & \code{FBV[0]} & $0\to1$ \\
      19 & \code{FBV[4]} & $0\to1$ \\
      \bottomrule
    \end{tabular}
    \par\smallskip
    \raggedright
    两处观察到同一receiver map；\par
    完整地址放讲稿。\par\medskip
    转换不是逐次命中记录；\par
    实际快路径命中率仍\par
    \textbf{UNKNOWN}。
  \end{column}
\end{columns}
\medskip
\textbf{待验证解释：}QJS重复own查找；V8可能复用稳定字段反馈。\par
现有证据尚未量出各步骤时间。
\end{frame}

% -------------------- Main 6/6 --------------------
\PageMeta{S1/S2正式汇总；S3 perturbation.json、site\_samples.csv、N1计数。}{\DiagNote}
\begin{frame}[t]{}
\small
\textbf{结论}
\begin{itemize}
  \setlength{\itemsep}{2pt}
  \item 操作不同，快慢方向可能不同；不把小程序拼成总体引擎排名。
  \item 方法案例中，多一层原型不足以解释跨引擎差距。
  \item 原Richards的普通own字段值得深入；已选出具体state位置。
\end{itemize}
\smallskip
\textbf{两个直接限制（诊断数据，不是正式性能结果）}
\begin{itemize}
  \setlength{\itemsep}{3pt}
  \item 无信号N512小对照：冻结版median1718ms，诊断开启版2223ms，\par
        高约\textbf{29\%}。不同批次点估计，不是精确插桩成本。
  \item 两次state读取每Run为10671/10115次，查找区域样本却为\textbf{85/9}。\par
        请求1ms、实际约10ms；采样与插桩影响分布，不能当精确耗时排名。
\end{itemize}
\smallskip
\textbf{问题}\par
先确认V8该位置实际快路径，还是先降低采样扰动、复核位置分布？
\end{frame}

\appendix
% -------------------- Appendix A --------------------
\PageMeta{S3：QuickJS完整目标函数dump与decoded.json；见notes.md。}{\DiagNote}
\begin{frame}[fragile,t]{附录A：QuickJS 完整最终指令序列}
\small 16条/29B；局部0=this，栈深2。保留全部指令与操作数，去掉源码注释。
\begin{lstlisting}
 0 push_this
 1 put_loc0 0: this
 2 get_loc0 0: this
 3 get_field state
 8 get_var 0: STATE_HELD
11 and
12 push_0 0
13 neq
14 dup
15 if_true8 28
17 drop
18 get_loc0 0: this
19 get_field state
24 get_var 1: STATE_SUSPENDED
27 eq
28 return
\end{lstlisting}
\end{frame}

% -------------------- Appendix B --------------------
\PageMeta{S3：同进程V8完整目标函数dump；IC地址映射见notes.md。}{\DiagNote}
\begin{frame}[fragile,t]{附录B：Ignition 完整最终指令序列}
\small 14条/31B；虚拟寄存器1、frame 8B。去掉地址/编码/调试列，指令与操作数完整。
\begin{lstlisting}
 0 GetNamedProperty <this>, [0:"state"], FBV[0]
 4 Star0
 5 LdaGlobal [1:"STATE_HELD"], FBV[2]
 8 BitwiseAnd r0, EmbeddedFeedback[None]
11 Star0
12 LdaZero
13 TestEqual r0, EmbeddedFeedback[None]
16 LogicalNot
17 JumpIfTrue [13] (0x1eef0100166e @ 30)
19 GetNamedProperty <this>, [0:"state"], FBV[4]
23 Star0
24 LdaGlobal [2:"STATE_SUSPENDED"], FBV[6]
27 TestEqual r0, EmbeddedFeedback[None]
30 Return
\end{lstlisting}
\small QJS 16条/29B，V8 14条/31B；条数、编码长度不是执行时间。
\end{frame}

\end{document}
